\QHeading{第1题　分一分。（填序号。）}

\par\noindent\makebox[\linewidth][c]{
\begin{tikzpicture}[x=1bp,y=-1bp,color=black,draw=black,text=black]
\begin{scope}\clip (56.0,108.0) rectangle (509.0,407.33);
\hdhead{41}{多边形的认识——信息窗1}
\hdsection{1}{129}{分一分。（填序号。）}
\begin{scope}[scale=.6666667,draw=black,text=black]\draw[hdoutline] (168.000,241.000)--(213.000,193.000)--(299.000,241.000)--cycle;\end{scope}
\begin{scope}[scale=.6666667,draw=black,text=black]\draw[hdoutline] (423.000,237.000)--(481.500,178.500)--(540.000,237.000)--cycle;\end{scope}
\begin{scope}[scale=.6666667,draw=black,text=black]\draw[hdoutline] (666.000,236.000)--(705.000,168.450)--(744.000,236.000)--cycle;\end{scope}
\begin{scope}[scale=.6666667,draw=black,text=black]\draw[hdoutline] (186.000,392.000)--(222.000,286.000)--(258.000,392.000)--cycle;\end{scope}
\begin{scope}[scale=.6666667,draw=black,text=black]\draw[hdoutline] (434.000,390.000)--(478.000,313.790)--(522.000,390.000)--cycle;\end{scope}
\begin{scope}[scale=.6666667,draw=black,text=black]\draw[hdoutline] (649.000,315.000)--(695.000,389.000)--(754.000,389.000)--cycle;\end{scope}
\begin{scope}[scale=.6666667,draw=black,text=black]\draw[line width=.45pt] (475.843,184.157)--(481.500,189.814)--(487.157,184.157);\end{scope}
\begin{scope}[scale=.6666667,draw=black,text=black]\node[inner sep=0pt,font=\normalsize,text=black] at (227.000,264.000) {①};\end{scope}
\begin{scope}[scale=.6666667,draw=black,text=black]\node[inner sep=0pt,font=\normalsize,text=black] at (482.000,264.000) {②};\end{scope}
\begin{scope}[scale=.6666667,draw=black,text=black]\node[inner sep=0pt,font=\normalsize,text=black] at (704.000,264.000) {③};\end{scope}
\begin{scope}[scale=.6666667,draw=black,text=black]\node[inner sep=0pt,font=\normalsize,text=black] at (222.000,416.000) {④};\end{scope}
\begin{scope}[scale=.6666667,draw=black,text=black]\node[inner sep=0pt,font=\normalsize,text=black] at (478.000,415.000) {⑤};\end{scope}
\begin{scope}[scale=.6666667,draw=black,text=black]\node[inner sep=0pt,font=\normalsize,text=black] at (703.000,412.000) {⑥};\end{scope}
\q{230}{438}{200}{13}{等腰三角形}
\q{572}{438}{200}{13}{等边三角形}
\begin{scope}[scale=.6666667,draw=black,text=black]\draw[rounded corners=12,dashed,line width=.6pt] (198,478) rectangle (383,583) (537,478) rectangle (718,583);\end{scope}
\hdsection{2}{614}{在下面的点子图中分别画一个直角三角形、一个锐角三角形和一个钝角三角形。}
\begin{scope}[scale=.6666667,draw=black,text=black]\foreach \i in {0,...,40}{\foreach \j in {0,...,13}{\fill[black] ({110+\i*17.5},{705+\j*17.5}) circle (1.15);}}\end{scope}
\hdsection{3}{976}{分别画出下面三角形底边上的高。}
\begin{scope}[scale=.6666667,draw=black,text=black]\draw[hdoutline] (170.000,1041.000)--(170.000,1121.000)--(286.000,1121.000)--cycle;\end{scope}
\begin{scope}[scale=.6666667,draw=black,text=black]\draw[hdoutline] (386.000,1121.000)--(510.000,1121.000)--(447.000,1052.000)--cycle;\end{scope}
\begin{scope}[scale=.6666667,draw=black,text=black]\draw[hdoutline] (597.000,1102.000)--(754.000,1039.000)--(686.000,1102.000)--cycle;\end{scope}
\begin{scope}[scale=.6666667,draw=black,text=black]\draw[line width=.45pt] (170.000,1112.000)--(179.000,1112.000)--(179.000,1121.000);\end{scope}
\begin{scope}[scale=.6666667,draw=black,text=black]\node[inner sep=0pt,font=\normalsize,text=black] at (243.000,1066.000) {底};\end{scope}
\begin{scope}[scale=.6666667,draw=black,text=black]\node[inner sep=0pt,font=\normalsize,text=black] at (404.000,1070.000) {底};\end{scope}
\begin{scope}[scale=.6666667,draw=black,text=black]\node[inner sep=0pt,font=\normalsize,text=black] at (643.000,1121.000) {底};\end{scope}

\begin{scope}[color=answerred,draw=answerred,text=answerred]\ans{151.33}{338.00}{96.00}{11}{②③④⑤}
\ans{382.67}{338.00}{73.33}{11}{③⑤}
\ans{136.67}{368.00}{106.00}{8}{等边三角形也是等腰三角形。}
\draw[answerred,line width=.8pt] (96.667,598.667) -- (96.667,505.333) -- (190.000,598.667) -- cycle;
\draw[answerred,line width=.7pt] (96.667,591.333) -- (104.000,591.333);
\draw[answerred,line width=.7pt] (104.000,591.333) -- (104.000,598.667);
\draw[answerred,line width=.8pt] (248.667,598.667) -- (342.000,598.667) -- (295.333,494.000) -- cycle;
\draw[answerred,line width=.8pt] (400.667,586.667) -- (517.333,586.667) -- (424.000,552.000) -- cycle;
\ans{104.67}{606.00}{86.67}{9}{直角三角形}
\ans{264.00}{606.00}{86.67}{9}{锐角三角形}
\ans{430.67}{602.00}{86.67}{9}{钝角三角形}
\draw[answerred,line width=.7pt,dashed] (113.333,747.333) -- (138.259,711.190);
\draw[answerred,line width=.7pt] (142.650,714.218) -- (139.622,718.609);
\draw[answerred,line width=.7pt] (139.622,718.609) -- (135.232,715.581);
\ans{129.80}{723.93}{20.00}{9}{高}
\draw[answerred,line width=.7pt,dashed] (340.000,747.333) -- (293.599,706.312);
\draw[answerred,line width=.7pt] (290.066,710.308) -- (294.062,713.840);
\draw[answerred,line width=.7pt] (294.062,713.840) -- (297.594,709.844);
\ans{320.80}{721.49}{20.00}{9}{高}
\draw[answerred,line width=.7pt,dashed] (457.333,734.667) -- (510.667,734.667);
\draw[answerred,line width=.7pt,dashed] (502.667,692.667) -- (502.667,734.667);
\draw[answerred,line width=.7pt] (497.333,734.667) -- (497.333,729.333);
\draw[answerred,line width=.7pt] (497.333,729.333) -- (502.667,729.333);
\ans{506.67}{708.33}{20.00}{9}{高}
\end{scope}
\end{scope}\pgfresetboundingbox\path[use as bounding box] (56.0,108.0) rectangle (509.0,407.33);
\end{tikzpicture}}\par
\par\noindent




\par

\par\noindent

\QHeading{第2题　在下面的点子图中分别画一个直角三角形、一个锐角三角形和一个钝角三角形。}

\par\noindent\makebox[\linewidth][c]{
\begin{tikzpicture}[x=1bp,y=-1bp,color=black,draw=black,text=black]
\begin{scope}\clip (56.0,433.0) rectangle (546.0,648.67);
\hdhead{41}{多边形的认识——信息窗1}
\hdsection{1}{129}{分一分。（填序号。）}
\begin{scope}[scale=.6666667,draw=black,text=black]\draw[hdoutline] (168.000,241.000)--(213.000,193.000)--(299.000,241.000)--cycle;\end{scope}
\begin{scope}[scale=.6666667,draw=black,text=black]\draw[hdoutline] (423.000,237.000)--(481.500,178.500)--(540.000,237.000)--cycle;\end{scope}
\begin{scope}[scale=.6666667,draw=black,text=black]\draw[hdoutline] (666.000,236.000)--(705.000,168.450)--(744.000,236.000)--cycle;\end{scope}
\begin{scope}[scale=.6666667,draw=black,text=black]\draw[hdoutline] (186.000,392.000)--(222.000,286.000)--(258.000,392.000)--cycle;\end{scope}
\begin{scope}[scale=.6666667,draw=black,text=black]\draw[hdoutline] (434.000,390.000)--(478.000,313.790)--(522.000,390.000)--cycle;\end{scope}
\begin{scope}[scale=.6666667,draw=black,text=black]\draw[hdoutline] (649.000,315.000)--(695.000,389.000)--(754.000,389.000)--cycle;\end{scope}
\begin{scope}[scale=.6666667,draw=black,text=black]\draw[line width=.45pt] (475.843,184.157)--(481.500,189.814)--(487.157,184.157);\end{scope}
\begin{scope}[scale=.6666667,draw=black,text=black]\node[inner sep=0pt,font=\normalsize,text=black] at (227.000,264.000) {①};\end{scope}
\begin{scope}[scale=.6666667,draw=black,text=black]\node[inner sep=0pt,font=\normalsize,text=black] at (482.000,264.000) {②};\end{scope}
\begin{scope}[scale=.6666667,draw=black,text=black]\node[inner sep=0pt,font=\normalsize,text=black] at (704.000,264.000) {③};\end{scope}
\begin{scope}[scale=.6666667,draw=black,text=black]\node[inner sep=0pt,font=\normalsize,text=black] at (222.000,416.000) {④};\end{scope}
\begin{scope}[scale=.6666667,draw=black,text=black]\node[inner sep=0pt,font=\normalsize,text=black] at (478.000,415.000) {⑤};\end{scope}
\begin{scope}[scale=.6666667,draw=black,text=black]\node[inner sep=0pt,font=\normalsize,text=black] at (703.000,412.000) {⑥};\end{scope}
\q{230}{438}{200}{13}{等腰三角形}
\q{572}{438}{200}{13}{等边三角形}
\begin{scope}[scale=.6666667,draw=black,text=black]\draw[rounded corners=12,dashed,line width=.6pt] (198,478) rectangle (383,583) (537,478) rectangle (718,583);\end{scope}
\hdsection{2}{614}{在下面的点子图中分别画一个直角三角形、一个锐角三角形和一个钝角三角形。}
\begin{scope}[scale=.6666667,draw=black,text=black]\foreach \i in {0,...,40}{\foreach \j in {0,...,13}{\fill[black] ({110+\i*17.5},{705+\j*17.5}) circle (1.15);}}\end{scope}
\hdsection{3}{976}{分别画出下面三角形底边上的高。}
\begin{scope}[scale=.6666667,draw=black,text=black]\draw[hdoutline] (170.000,1041.000)--(170.000,1121.000)--(286.000,1121.000)--cycle;\end{scope}
\begin{scope}[scale=.6666667,draw=black,text=black]\draw[hdoutline] (386.000,1121.000)--(510.000,1121.000)--(447.000,1052.000)--cycle;\end{scope}
\begin{scope}[scale=.6666667,draw=black,text=black]\draw[hdoutline] (597.000,1102.000)--(754.000,1039.000)--(686.000,1102.000)--cycle;\end{scope}
\begin{scope}[scale=.6666667,draw=black,text=black]\draw[line width=.45pt] (170.000,1112.000)--(179.000,1112.000)--(179.000,1121.000);\end{scope}
\begin{scope}[scale=.6666667,draw=black,text=black]\node[inner sep=0pt,font=\normalsize,text=black] at (243.000,1066.000) {底};\end{scope}
\begin{scope}[scale=.6666667,draw=black,text=black]\node[inner sep=0pt,font=\normalsize,text=black] at (404.000,1070.000) {底};\end{scope}
\begin{scope}[scale=.6666667,draw=black,text=black]\node[inner sep=0pt,font=\normalsize,text=black] at (643.000,1121.000) {底};\end{scope}

\begin{scope}[color=answerred,draw=answerred,text=answerred]\ans{151.33}{338.00}{96.00}{11}{②③④⑤}
\ans{382.67}{338.00}{73.33}{11}{③⑤}
\ans{136.67}{368.00}{106.00}{8}{等边三角形也是等腰三角形。}
\draw[answerred,line width=.8pt] (96.667,598.667) -- (96.667,505.333) -- (190.000,598.667) -- cycle;
\draw[answerred,line width=.7pt] (96.667,591.333) -- (104.000,591.333);
\draw[answerred,line width=.7pt] (104.000,591.333) -- (104.000,598.667);
\draw[answerred,line width=.8pt] (248.667,598.667) -- (342.000,598.667) -- (295.333,494.000) -- cycle;
\draw[answerred,line width=.8pt] (400.667,586.667) -- (517.333,586.667) -- (424.000,552.000) -- cycle;
\ans{104.67}{606.00}{86.67}{9}{直角三角形}
\ans{264.00}{606.00}{86.67}{9}{锐角三角形}
\ans{430.67}{602.00}{86.67}{9}{钝角三角形}
\draw[answerred,line width=.7pt,dashed] (113.333,747.333) -- (138.259,711.190);
\draw[answerred,line width=.7pt] (142.650,714.218) -- (139.622,718.609);
\draw[answerred,line width=.7pt] (139.622,718.609) -- (135.232,715.581);
\ans{129.80}{723.93}{20.00}{9}{高}
\draw[answerred,line width=.7pt,dashed] (340.000,747.333) -- (293.599,706.312);
\draw[answerred,line width=.7pt] (290.066,710.308) -- (294.062,713.840);
\draw[answerred,line width=.7pt] (294.062,713.840) -- (297.594,709.844);
\ans{320.80}{721.49}{20.00}{9}{高}
\draw[answerred,line width=.7pt,dashed] (457.333,734.667) -- (510.667,734.667);
\draw[answerred,line width=.7pt,dashed] (502.667,692.667) -- (502.667,734.667);
\draw[answerred,line width=.7pt] (497.333,734.667) -- (497.333,729.333);
\draw[answerred,line width=.7pt] (497.333,729.333) -- (502.667,729.333);
\ans{506.67}{708.33}{20.00}{9}{高}
\end{scope}
\end{scope}\pgfresetboundingbox\path[use as bounding box] (56.0,433.0) rectangle (546.0,648.67);
\end{tikzpicture}}\par



\par





\par\noindent

\QHeading{第2题　直角、锐角、钝角三角形的另一些示例}

\Answer{以下三个图形与前面的示例不同，但仍分别满足直角、锐角、钝角条件。作图答案不唯一。}
\par\smallskip\noindent\begin{adjustbox}{max width=\linewidth,center}
\begin{tikzpicture}[x=.55cm,y=.55cm,font=\normalsize]
\draw[answerred,thick] (0,0)--(3,0)--(0,2)--cycle;
\draw[answerred] (0,.25)--(.25,.25)--(.25,0);
\node[below] at (1.5,0){直角三角形};
\begin{scope}[xshift=2.7cm]\draw[answerred,thick] (0,0)--(3,0)--(1.5,2.5)--cycle;\node[below] at (1.5,0){锐角三角形};\end{scope}
\begin{scope}[xshift=5.4cm]\draw[answerred,thick] (0,0)--(3,0)--(1,.5)--cycle;\node[below] at (1.5,0){钝角三角形};\end{scope}
\end{tikzpicture}
\end{adjustbox}\par\smallskip

\QHeading{第3题　分别画出下面三角形底边上的高。}

\par\noindent\makebox[\linewidth][c]{
\begin{tikzpicture}[x=1bp,y=-1bp,color=black,draw=black,text=black]
\begin{scope}\clip (108.0,687.0) rectangle (520.0,758.0);
\hdhead{41}{多边形的认识——信息窗1}
\hdsection{1}{129}{分一分。（填序号。）}
\begin{scope}[scale=.6666667,draw=black,text=black]\draw[hdoutline] (168.000,241.000)--(213.000,193.000)--(299.000,241.000)--cycle;\end{scope}
\begin{scope}[scale=.6666667,draw=black,text=black]\draw[hdoutline] (423.000,237.000)--(481.500,178.500)--(540.000,237.000)--cycle;\end{scope}
\begin{scope}[scale=.6666667,draw=black,text=black]\draw[hdoutline] (666.000,236.000)--(705.000,168.450)--(744.000,236.000)--cycle;\end{scope}
\begin{scope}[scale=.6666667,draw=black,text=black]\draw[hdoutline] (186.000,392.000)--(222.000,286.000)--(258.000,392.000)--cycle;\end{scope}
\begin{scope}[scale=.6666667,draw=black,text=black]\draw[hdoutline] (434.000,390.000)--(478.000,313.790)--(522.000,390.000)--cycle;\end{scope}
\begin{scope}[scale=.6666667,draw=black,text=black]\draw[hdoutline] (649.000,315.000)--(695.000,389.000)--(754.000,389.000)--cycle;\end{scope}
\begin{scope}[scale=.6666667,draw=black,text=black]\draw[line width=.45pt] (475.843,184.157)--(481.500,189.814)--(487.157,184.157);\end{scope}
\begin{scope}[scale=.6666667,draw=black,text=black]\node[inner sep=0pt,font=\normalsize,text=black] at (227.000,264.000) {①};\end{scope}
\begin{scope}[scale=.6666667,draw=black,text=black]\node[inner sep=0pt,font=\normalsize,text=black] at (482.000,264.000) {②};\end{scope}
\begin{scope}[scale=.6666667,draw=black,text=black]\node[inner sep=0pt,font=\normalsize,text=black] at (704.000,264.000) {③};\end{scope}
\begin{scope}[scale=.6666667,draw=black,text=black]\node[inner sep=0pt,font=\normalsize,text=black] at (222.000,416.000) {④};\end{scope}
\begin{scope}[scale=.6666667,draw=black,text=black]\node[inner sep=0pt,font=\normalsize,text=black] at (478.000,415.000) {⑤};\end{scope}
\begin{scope}[scale=.6666667,draw=black,text=black]\node[inner sep=0pt,font=\normalsize,text=black] at (703.000,412.000) {⑥};\end{scope}
\q{230}{438}{200}{13}{等腰三角形}
\q{572}{438}{200}{13}{等边三角形}
\begin{scope}[scale=.6666667,draw=black,text=black]\draw[rounded corners=12,dashed,line width=.6pt] (198,478) rectangle (383,583) (537,478) rectangle (718,583);\end{scope}
\hdsection{2}{614}{在下面的点子图中分别画一个直角三角形、一个锐角三角形和一个钝角三角形。}
\begin{scope}[scale=.6666667,draw=black,text=black]\foreach \i in {0,...,40}{\foreach \j in {0,...,13}{\fill[black] ({110+\i*17.5},{705+\j*17.5}) circle (1.15);}}\end{scope}
\hdsection{3}{976}{分别画出下面三角形底边上的高。}
\begin{scope}[scale=.6666667,draw=black,text=black]\draw[hdoutline] (170.000,1041.000)--(170.000,1121.000)--(286.000,1121.000)--cycle;\end{scope}
\begin{scope}[scale=.6666667,draw=black,text=black]\draw[hdoutline] (386.000,1121.000)--(510.000,1121.000)--(447.000,1052.000)--cycle;\end{scope}
\begin{scope}[scale=.6666667,draw=black,text=black]\draw[hdoutline] (597.000,1102.000)--(754.000,1039.000)--(686.000,1102.000)--cycle;\end{scope}
\begin{scope}[scale=.6666667,draw=black,text=black]\draw[line width=.45pt] (170.000,1112.000)--(179.000,1112.000)--(179.000,1121.000);\end{scope}
\begin{scope}[scale=.6666667,draw=black,text=black]\node[inner sep=0pt,font=\normalsize,text=black] at (243.000,1066.000) {底};\end{scope}
\begin{scope}[scale=.6666667,draw=black,text=black]\node[inner sep=0pt,font=\normalsize,text=black] at (404.000,1070.000) {底};\end{scope}
\begin{scope}[scale=.6666667,draw=black,text=black]\node[inner sep=0pt,font=\normalsize,text=black] at (643.000,1121.000) {底};\end{scope}

\begin{scope}[color=answerred,draw=answerred,text=answerred]\ans{151.33}{338.00}{96.00}{11}{②③④⑤}
\ans{382.67}{338.00}{73.33}{11}{③⑤}
\ans{136.67}{368.00}{106.00}{8}{等边三角形也是等腰三角形。}
\draw[answerred,line width=.8pt] (96.667,598.667) -- (96.667,505.333) -- (190.000,598.667) -- cycle;
\draw[answerred,line width=.7pt] (96.667,591.333) -- (104.000,591.333);
\draw[answerred,line width=.7pt] (104.000,591.333) -- (104.000,598.667);
\draw[answerred,line width=.8pt] (248.667,598.667) -- (342.000,598.667) -- (295.333,494.000) -- cycle;
\draw[answerred,line width=.8pt] (400.667,586.667) -- (517.333,586.667) -- (424.000,552.000) -- cycle;
\ans{104.67}{606.00}{86.67}{9}{直角三角形}
\ans{264.00}{606.00}{86.67}{9}{锐角三角形}
\ans{430.67}{602.00}{86.67}{9}{钝角三角形}
\draw[answerred,line width=.7pt,dashed] (113.333,747.333) -- (138.259,711.190);
\draw[answerred,line width=.7pt] (142.650,714.218) -- (139.622,718.609);
\draw[answerred,line width=.7pt] (139.622,718.609) -- (135.232,715.581);
\ans{129.80}{723.93}{20.00}{9}{高}
\draw[answerred,line width=.7pt,dashed] (340.000,747.333) -- (293.599,706.312);
\draw[answerred,line width=.7pt] (290.066,710.308) -- (294.062,713.840);
\draw[answerred,line width=.7pt] (294.062,713.840) -- (297.594,709.844);
\ans{320.80}{721.49}{20.00}{9}{高}
\draw[answerred,line width=.7pt,dashed] (457.333,734.667) -- (510.667,734.667);
\draw[answerred,line width=.7pt,dashed] (502.667,692.667) -- (502.667,734.667);
\draw[answerred,line width=.7pt] (497.333,734.667) -- (497.333,729.333);
\draw[answerred,line width=.7pt] (497.333,729.333) -- (502.667,729.333);
\ans{506.67}{708.33}{20.00}{9}{高}
\end{scope}
\end{scope}\pgfresetboundingbox\path[use as bounding box] (108.0,687.0) rectangle (520.0,758.0);
\end{tikzpicture}}\par



\par







\QHeading{第3题　另一个作图检查方法}

\Answer{先检查是否用虚线画高、标“高”和直角符号，再检查高是否经过与指定底相对的顶点，再用三角板检查是否垂直于底边所在直线。若垂足在底边外，底边的延长线用虚线；不能为了把高画到图形内部而改变指定的底。}
